\documentclass[10pt,a4paper]{amsart}
\usepackage[margin=19mm]{geometry}
\usepackage{amsmath,amssymb,mathtools}
\usepackage[scr=rsfs,cal=euler]{mathalfa}
\usepackage{xfrac}
\usepackage[unicode,hidelinks,bookmarks=false]{hyperref}
\newcommand{\ie}{i.e.\ }
\newcommand{\cf}{cf.\ }
\newcommand{\resp}{respectively\ }
\newcommand{\Subsection}{\subsection}
\providecommand{\nobreakdash}{\nobreak\hbox{-}}
\setlength{\emergencystretch}{2em}
\multlinegap0pt
\usepackage{bm}
\usepackage{bbm}
\usepackage{dsfont}
\usepackage{xspace}
\usepackage{cancel}
\usepackage{mathtools}
\usepackage{amsthm}
\makeatletter
\@ifundefined{lemma}{\newtheorem{lemma}{Lemma}}{}
\@ifundefined{notation}{\newtheorem{notation}{Notation}}{}
\@ifundefined{theorem}{\newtheorem{theorem}{Theorem}}{}
\makeatother
\makeatletter
\newcommand{\E}{\mathbb{E}}
\newcommand{\F}{\mathbb{F}}
\expandafter\ifx\csname clproof\endcsname\relax
\newenvironment{clproof}[1][\myproofname]{\begin{proof}[#1]\renewcommand*{\qedsymbol}{\(\blacktriangle\)}}{\end{proof}}
\fi
\newcommand{\mA}{\mathcal{A}}
\newcommand{\mB}{\mathcal{B}}
\newcommand{\mS}{\mathcal{S}}
\newcommand{\tmax}{T_{\max}}
\makeatother
\title[A Combinatorial Perspective on Random Access Efficiency for ]{A Combinatorial Perspective on Random Access Efficiency for DNA Storage}
\author{Anina Gruica, Daniella Bar-Lev, Alberto Ravagnani, Eitan Yaakobi}
\date{Source: arXiv:2401.15722v3 (CC BY 4.0)}
\begin{document}
\maketitle

\noindent\emph{Source and licence.} A Combinatorial Perspective on Random Access Efficiency for DNA Storage. Version arXiv:2401.15722v3 (CC BY 4.0), distributed under the Creative Commons Attribution 4.0 International licence, as recorded in the arXiv licence metadata for this exact version and shown on the abstract page https://arxiv.org/abs/2401.15722v3. Full rights notice: licenses/arxiv-2401.15722-CC-BY-4.0.txt.\par\smallskip

\noindent\textit{Editorial scope.} Counted unit: the Theorem whose source label is \texttt{prop:extmds} in the article (source0.tex lines 1476--1481, source label \texttt{prop:extmds}), delivered together with the article's own complete original proof of it (source0.tex lines 1482--1540, one \texttt{proof} environment of 59 lines). No mathematics is written, repaired or re-derived: statement and proof are the article's own text line for line. Required context retained uncounted: lem:fi (lines 346--351); not:extmds (lines 1464--1466). Every definition, display and label of those lines is retained unchanged. Omitted from this package and not redistributed: 1--345, 352--1463, 1467--1475, 1541--2480. Nothing inside a retained excerpt is omitted or altered: every retained body is delivered exactly as the article's lines are written, so no omission receipt of any kind accompanies this package. Citations: the retained excerpts carry no citation command at all, so no bibliography entry of the article is transcribed and no reference is redistributed. Reference adaptation: none was needed and none was made; every reference inside the delivered unit resolves inside it, so no unresolved reference or citation is produced. Printed slips kept literal: no printed word, symbol, hypothesis or number of the counted statement or of its proof was altered. Layout and class: the article's own class is \texttt{amsart} with packages including amssymb, amsmath, xcolor, float, latexsym, amsthm, empheq, bm, booktabs, subcaption, tikz, pgfplots, cite, balance; the wrapper is the lane's pinned \texttt{amsart} editorial wrapper at 10pt on A4, which restates the theorem-like environments the retained text needs and adds the article's own macro definitions that the retained text needs (\texttt{\textbackslash E}, \texttt{\textbackslash F}, \texttt{\textbackslash mA}, \texttt{\textbackslash mB}, \texttt{\textbackslash mS}, \texttt{\textbackslash tmax}), with extra wrapper packages (bm, bbm, dsfont, xspace, cancel, mathtools). The delivered numbering is the unit's own. Assets: no retained span contains a figure environment, a graphics-inclusion command, a PDF-page inclusion, a table or a TikZ picture; no image file is copied into this package, the package declares no asset dependency and no image file is redistributed with it. Licence: version arXiv:2401.15722v3 of this article is distributed under the Creative Commons Attribution 4.0 International licence, as recorded in the arXiv licence metadata for this exact version and shown on the abstract page https://arxiv.org/abs/2401.15722v3. A full rights notice is delivered at licenses/arxiv-2401.15722-CC-BY-4.0.txt. Characters: every retained excerpt is pure ASCII, so the delivered engine, XeLaTeX with its default Latin Modern text font, reports no missing character.
\begin{lemma} \label{lem:fi}
For $G\in \F_q^{k \times n}$ and for all $i\in[k]$ we have
\begin{align*}
    \E[\tau_i(G)] =  n H_n - \sum_{s=1}^{n-1} \displaystyle \frac{\alpha_i(s)}{\binom{n-1}{s}}.
\end{align*}
\end{lemma}

\begin{notation} \label{not:extmds}
Let $G= (I_k \mid R) \in \F_q^{k \times n}$ be a systematic generator matrix of a code. For $x \ge 1$, we let $G^x = (I_k \mid I_k \mid \dots \mid I_k \mid R) \in \F_q^{k \times N}$ (where $N=xk+n-k$) be the matrix obtained by appending additional $x-1$ identity matrices to $G$. 
\end{notation}

\begin{theorem} \label{prop:extmds}
Let $G= (I_k \mid R) \in \F_q^{k \times n}$ be a systematic generator matrix of an MDS code and let $N=xk+n-k$. We have
\begin{align*}
\tmax(G^x) = 1+\sum_{s=1}^{N-1}\frac{\binom{N-x}{s}}{\binom{N-1}{s}} -\sum_{s=k}^{N-1}\sum_{a=0}^{k-1}\frac{\binom{k-1}{a}}{\binom{N-1}{s}} \sum_{m=0}^{s-k}\binom{n-k}{s-a-m}\sum_{t=0}^a(-1)^t\binom{a}{t}\binom{(a-t)x}{m+a}.
\end{align*}
\end{theorem}

\begin{proof}
In order to give an expression for $\E[\tau_i(G^x)]$, by (the natural analogue of) Lemma~\ref{lem:fi} it suffices to compute $\widetilde{\alpha}_i(s)$
for the code generated by $G^x$ and for $1\leq i\leq k$. Since $G$ is a systematic generator matrix of an MDS code, the only possible ways to recover $e_i$ is by either sampling one of the columns corresponding to $e_i$ itself, or by sampling at least~$k$ distinct columns of $G^x$. If $1\le s \le k-1$, it is not hard to see that 
\begin{align*}
        \widetilde{\alpha}_i(s) = \binom{N}{s}-\binom{N-x}{s}.
\end{align*}
We then focus on the case
$k\le s \le N-1$.
We can recover $e_i$ from $s$ columns of $G^x$ if~$e_i$ is one of them. The only other way to recover $e_i$ from~$s$ columns of $G^x$ is by sampling at least~$k$ distinct columns. Therefore, we need to find the size of the set
    $\{S \subseteq [N] : |S|=s, \, e_i \notin \{g_j : j \in S\}, \, |\{ g_j : j \in S\}| \ge k\}$,
    which we denote by $\mS$ in the sequel.
    
    Let $\mA=[xk]\setminus\{i,2i,\dots,xi\}$ and $\mB=\{xk+1,\dots,N\}$, so that we have $g_j=e_j$ for all $j \in \mA$ and the columns indexed by $\mB$ correspond to the columns of $R$, which is the redundancy part of $G=(I_k \mid R)$ as in Notation~\ref{not:extmds}.
    We have that $|\mS|$ is
    \begin{align} \label{eq:abab}
         \sum_{a=0}^{k-1} & |\{(A,B) \in 2^{\mA} \times 2^{\mB}  : |\{ g_j : j \in A\}|=a, |B| = s-|A|\}| \nonumber \\
        &=\sum_{a=0}^{k-1}\sum_{m=0}^{s-k}\sum_{\substack{A \subseteq \mA \\ |A|=a+m \\ |\{ g_j : j \in A\}|=a}}|\{B \subseteq \mB :   |B| = s-|A|\}| \nonumber \\
        &= \sum_{a=0}^{k-1}\sum_{m=0}^{s-k}\sum_{\substack{A \subseteq \mA \\ |A|=a+m \\ |\{ g_j : j \in A\}|=a}} \binom{N-xk}{s-a-m}.
    \end{align}
    For any $a \in \{0, \dots, k-1\}$ and $m \in \{0,\dots,s-k\}$, we are left with computing the size of the set
    \begin{align*}
        |\{A \subseteq \mA : |\{ g_j : j \in A\}|=a, |A|=a+m\}|.
    \end{align*}
    Let $C \subseteq \{g_j : j \in \mA\}$ with $|C|=a$. Using the Inclusion-Exclusion Principle, one can verify that we have
    \begin{align*}
        &|\{A \subseteq \mA : \{ g_j : j \in A\}=C, |A|=a+m\}| \\
        &= \sum_{t=0}^a(-1)^t\binom{a}{t}\binom{(a-t)x}{m+a}.
    \end{align*}
    Summing over all $C \subseteq \{g_j : j \in \mA\}$ with $|C|=a$ together with~\eqref{eq:abab} gives that $|\mS|$ equals
    \begin{align*}
      \sum_{a=0}^{k-1}\binom{k-1}{a} \sum_{m=0}^{s-k}\binom{N-xk}{s-a-m}\sum_{t=0}^a(-1)^t\binom{a}{t}\binom{(a-t)x}{m+a}.
    \end{align*}
By combining all of this
with Lemma~\ref{lem:fi} we finally obtain
\begin{align*}
\E[\tau_i(G^x)] = NH_N-\sum_{s=1}^{N-1}\frac{\binom{N}{s}-\binom{N-x}{s}}{\binom{N-1}{s}}-\sum_{s=k}^{N-1}\frac{|\mS|}{\binom{N-1}{s}},
\end{align*}
which does not depend on the coordinate $i$ and therefore, after simplifying, gives the statement of the proposition.
% In order to give an expression for $\E[\tau_i(G^x)]$, by (the natural analogue of) Lemma~\ref{lem:fi} it suffices to compute $\widetilde{\alpha}_i(s)$
% for the code generated by $G^x$. Since $G$ is a generator matrix of an MDS code, the only possible ways to recover the $i$-th information strand is by either sampling the strand itself, or by sampling at least~$k$ distinct columns of $G^x$. If $1\le s \le k-1$, it is not hard to see that 
% \begin{align*}
%         \widetilde{\alpha}_i(s) = \binom{N}{s}-\binom{N-x}{s}.
%     \end{align*}
% We then focus on the case
% $k\le s \le N-1$.
% We can recover $e_i$ from $s$ columns of $G^x$ if~$e_i$ is one of them. The only other way to recover $e_i$ from~$s$ columns of $G^x$ is by sampling at least~$k$ distinct columns. Therefore we need to find the size the set
%     $\{S \subseteq [N] : |S|=s, \, e_i \notin \{g_j : j \in S\}, \, |\{ g_j : j \in S\}| \ge k\}$,
%     which we denote by $\mS$ in the sequel. Using the Inclusion-Exclusion Principle, together with some technical rewriting, we obtain that $|\mS|$ equals
%     \begin{align*}
%       \sum_{a=0}^{k-1}\binom{k-1}{a} \sum_{m=0}^{s-k}\binom{N-xk}{s-a-m}\sum_{t=0}^a(-1)^t\binom{a}{t}\binom{(a-t)x}{m+a};
%     \end{align*}
% see~\cite{gruica2024reducing} for a more detailed computation for the value of $|\mS|$.
% By combining all of this
% with Lemma~\ref{lem:fi} we finally obtain
% \begin{align*}
% \E[\tau_i(G^x)] = NH_N-\sum_{s=1}^{N-1}\frac{\binom{N}{s}-\binom{N-x}{s}}{\binom{N-1}{s}}-\sum_{s=k}^{N-1}\frac{|\mS|}{\binom{N-1}{s}},
% \end{align*}
% which does not depend on the coordinate $i$ and therefore, after simplifying, gives the statement of the proposition.
\end{proof}

\end{document}
